\begin{abstract}
Worldwide image geolocation models such as GeoCLIP estimate GPS coordinates
from a single photograph with high accuracy, yet they provide no account of
which image evidence supports a given prediction. Post-hoc saliency methods
address this question only correlationally: they indicate where activation
concentrates, without establishing that the highlighted evidence is what the
model relies upon. This paper presents GeoProbe, an interactive system for
causal analysis of image geolocation models. Given an image region ---
specified manually, or obtained from open-vocabulary Grounding DINO
detections over a vocabulary of geographic cues such as \texttt{street
sign}, \texttt{utility pole} and \texttt{license plate} --- GeoProbe adds a
key-indexed bias to the pre-softmax attention logits of the CLIP vision
tower, independently per layer, thereby increasing or decreasing the
attention mass allocated to that region. The system executes GeoCLIP twice,
without and with the intervention, and reports both predictions together
with the geodesic distance separating them. The magnitude of this
displacement constitutes a direct causal measure of the region's
contribution to the prediction. We further report a search over attention
biases applied to early, middle and late layer groups on the Img2GPS3k
benchmark, quantifying the extent to which the geolocation output can be
manipulated through attention reweighting alone.
\TODO{add one sentence stating the principal quantitative result}

\keywords{Image Geolocation \and Explainability \and Attention Intervention
\and Vision--Language Models.}
\end{abstract}
